\chapter{Conclusiones y posibles ampliaciones}


\section{Principales aportaciones y limitaciones}

Tras el análisis crítico de los experimentos ejecutados, se presenta a continuación una síntesis de las principales conclusiones y aportaciones técnicas derivadas del diseño e implementación del sistema:

\begin{itemize}
    \item \textbf{Sinergia de los reordenamientos con las estructuras de almacenamiento:} La poda local actúa de forma dependiente con la estrategia de gestión de memoria. Manifestó una ganancia crítica al acoplarse al algoritmo orientado a bloques (\textit{neighborsStruct}), ya que la discriminación geométrica se especializa en agrupar puntos vecinos y reduce las bifurcaciones condicionales y las operaciones de inserción. Esto provoca que el reordenamiento cartesiano no produce resultados favorables al no poder acoplarse naturalmente a este sistema de inserciones en bloque, al contrario que el reordenamiento propuesto mediante el ángulo $\varphi$.
    
    \item \textbf{Especialización y dinámica dimensional de los kernels:} Se comprueba que la efectividad teórica del descarte está ligada a la morfología de la consulta. El modo polar maximiza su rendimiento en kernels esféricos (especialmente a radios bajos debido al estrechamiento de la sección angular), mientras que el modo cartesiano destaca en kernels cúbicos (particularmente a radios mayores, donde los cortes unidimensionales con las caras del nodo permiten una separación perfecta entre puntos vecinos y no vecinos).
    
    \item \textbf{Delimitación del espacio de hiperparámetros:} A pesar de la variabilidad y la naturaleza única de cada nube LiDAR, se consolida un umbral operativo óptimo fijando el \texttt{threshold} en 64 puntos, cota que evita introducir un \textit{overhead} ineficiente en la CPU ante subconjuntos minúsculos. Asimismo, se establece que valores superiores a \texttt{maxLeaf = 256} marcan el inicio de la degradación de la eficiencia, partiendo el espacio en octantes no lo suficientemente concretos. Aún así, esta implementación supone una ganancia en flexibilidad a la hora de diseñar el tamaño de los octantes del \textit{Octree}.
    
    \item \textbf{Inversión de rentabilidad en coberturas completas:} Se identifica una limitación física del algoritmo en las búsquedas completas. La ordenación global por curvas de Hilbert combinada con la asignación por bloques de OpenMP maximiza de forma nativa la reutilización de datos en la memoria caché de la CPU y el tiempo ahorrado con los descartes de puntos no llega a equilibrar la sobrecarga introducida.
\end{itemize}

A partir de estas premisas, se puede confirmar la mejora real introducida en este trabajo mediante la propuesta del reordenamiento basado en coordenadas esféricas integrado sobre el algoritmo \textit{neighborsStruct}. Para poner en el contexto de la literatura técnica esta implementación, y tal como se abordó en el estado del arte, se tomaron como referencia de rendimiento los resultados del Octree lineal publicados por Viñambres et al. \cite{Vinambres}. Dicha estructura base minimizaba los tiempos de ejecución en relación con las principales implementaciones abiertas de Octrees, tales como PCL, PicoTree, NanoFLANN y unibnOctree. 

Por consiguiente, la integración original de la técnica de poda a nivel de hoja desarrollada en este proyecto sobre la arquitectura de Viñambres et al. supone, hasta donde se conoce y bajo el entorno operativo adecuado (consultas individuales/aleatorias con kernels esféricos), la solución más eficiente en la actualidad para detectar y almacenar vecinos sobre nubes de puntos tridimensionales LiDAR.

\section{Vías de Mejora y Trabajo Futuro}

Por último, como una ampliación a futuro, se propone de desarrollo el diseño de un mecanismo que permita acoplar la búsqueda de rangos cartesiana directamente sobre la infraestructura orientada a datos contiguos de \textit{neighborsStruct}.

Esta integración tiene como objetivo desbloquear el gran potencial de la poda por ejes cartesianos sobre geometría cúbicas, visible en los porcentajes de puntos podados conseguidos, cuyo impacto no pudo transformarse en eficiencia temporal real. El motivo es la incompatibilidad en cuanto a estructura del algoritmo \textit{neighborsStruct} base con los múltiples vectores reordenados usados simultáneamente durante las consultas. En caso de encontrar una solución para esta integración, se elevaría drásticamente la rapidez temporal de las consultas sobre kernels cúbicos.

Teóricamente, esta eficiencia sería incluso de mayor magnitud que la encontrada para kernels esféricos mediante la poda angular. Esto es debido a que la efectividad geométrica del descarte cartesiano es sustancialmente más alta: tal y como se constató en las métricas de depuración, apenas existen puntos evaluados por el algoritmo que no terminen cumpliendo el criterio de vecindad. Esto implica que el sistema haría todavía menos operaciones de cómputo de vecinos por hoja. Esto constituye, sin duda, la línea de investigación más prometedora para la continuidad de este trabajo.
